% Einheit 1 von 4 – Mittlere Änderungsrate und Sekante (bank/ableitung-und-aenderungsrate/e1.jsonl, zusammenbau v0.1)
%% TODO Zeitmarke Einheit 1: Marken-Zeile nicht lesbar
%% TODO Zweigzeile Teil 1: was hier gelernt wird (Satz in Schülersprache aus dem Einheitstitel) – Einheit 1
\einheitenkopf[e1]{Einheit 1 von 4 · Mittlere Änderungsrate und Sekante}
\zweigzeile{Abitur GK}
\setcounter{aufgabe}{9}

%% TODO Titel: Ich-kann-Satz fehlt in der Bank (Platzhalter: Kettenname) – Nr. 10
%% TODO Anweisung: ein Satz über der ersten Teilaufgabe fehlt in der Bank – Nr. 10
\begin{aufgabe}{Mittlere Änderungsrate}
%% TODO \rechenplatz{n}: Schrittzahl des Grundfalls fehlt in der Bank (nur bei fester Schreibform, 2.3 b)
\begin{teile}
\teil „Wie stark hat der Wasserstand in den ersten fünf Stunden durchschnittlich je Stunde zugenommen?“ \\ \kreuz{Zeitraum – Differenzenquotient, Sekante} \\ \kreuz{Zeitpunkt – Ableitung, Tangente}
\teil Um 8 Uhr sind 120 Besucher im Zoo, um 11 Uhr 300. Wie groß ist die mittlere Änderungsrate der Besucherzahl in diesem Zeitraum? \leerfeld Besucher je Stunde
\teil Die Tabelle zeigt, wie viele Anmeldungen bis zur jeweiligen Uhrzeit eingegangen sind. Berechne die mittlere Anzahl der Anmeldungen je Stunde zwischen 7:45 Uhr und 9:15 Uhr \hfill (Abitur 2024 LK).

\sachtabelle{lccccc}{Uhrzeit & 7:00 & 7:45 & 8:30 & 9:15 & 10:00}{Anmeldungen & 210 & 480 & 890 & 1\,530 & 2\,100}
\leerfeld je Stunde
\teil Der Hormonspiegel eines Patienten wird durch $h(t) = 6t \cdot \mathrm{e}^{-0{,}05t} + 40$ beschrieben ($t$ in Tagen seit Behandlungsbeginn, $h(t)$ in Prozent des Sollwerts). Berechne die mittlere Änderungsrate des Hormonspiegels in den ersten fünf Tagen \hfill (Abitur 2019 GK). \leerfeld Prozentpunkte je Tag
\teil Der Wasserstand eines Bachs wird durch $h(t) = 0{,}5t^3 - 6t^2 + 20t + 30$ beschrieben ($t$ in Stunden seit Regenbeginn, $h(t)$ in cm). Berechne den Wert des Terms $\frac{1}{3} \cdot (h(5) - h(2))$ und deute ihn im Sachzusammenhang \hfill (Abitur 2025 LK).
\teil Ein Höhenprofil wird für $0 \le x \le 6$ durch $f(x) = (x + 2) \cdot \mathrm{e}^{-0{,}4x}$ beschrieben (1 LE = 1 km). Im Intervall $[1; 5]$ soll das Profil näherungsweise durch die Gerade durch die Punkte $(1 | f(1))$ und $(5 | f(5))$ ersetzt werden. Gib eine Gleichung dieser Geraden an \hfill (Abitur 2018 GK).
\teil Ein Höhenprofil wird durch $f(x) = (x + 2) \cdot \mathrm{e}^{-0{,}4x}$ beschrieben (1 LE = 1 km). Von einem zweiten Profil $h$ ist bekannt: $h(0) = 1{,}8$ und $h(5) = 0{,}6$. Vergleiche die Beträge der mittleren Steigungen von $f$ und $h$ im Intervall $[0; 5]$ \hfill (Abitur 2018 GK).
\teil Der Schienenverlauf einer Modelleisenbahn besteht aus vier gleichen Bögen von je $3{,}2$ m Länge und zwei Halbkreisen mit dem Radius $0{,}5$ m, die die Bögen verbinden. Ein Zug fährt im Durchschnitt $4$ m je Minute. Wie lange braucht er für eine volle Runde? \hfill (Abitur 2026 GK) \\ \leerfeld[min]
\teil Der Bestand eines Rohstoffs wird durch $k(t) = 0{,}1t^3 - 1{,}5t^2 + 100$ beschrieben ($t$ in Jahren, $k(t)$ in Mio. Tonnen). Die Gleichung $\frac{k(c) - k(0)}{c} = -5$ hat genau zwei Lösungen. Bestimme die kleinere Lösung und deute Gleichung und Lösung im Sachzusammenhang \hfill (Abitur 2026 GK).
\teil Der Graph zeigt das Höhenprofil eines Skateparks (1 LE = 1 m). Kriterium: Die Gerade durch Anfangs- und Endpunkt darf gegen die Waagerechte höchstens $45°$ geneigt sein. Zeichne die $45°$-Gerade durch den Anfangspunkt und die Sekante ein und beurteile ohne Rechnung, ob das Kriterium erfüllt ist \hfill (Abitur 2024 LK).

\begin{ksys}[xmin=0,xmax=4,ymin=0,ymax=4,ablesen]\funktion{0.1*\x^3}{f}\end{ksys}
\teil Die Kosten für die Produktion von $x$ Kubikmetern einer Flüssigkeit werden für $0 \le x \le 8$ durch $K(x) = x^3 - 9x^2 + 30x + 10$ beschrieben ($K(x)$ in 1\,000 Euro). Aussage: „Je größer die Produktionsmenge, desto höher sind die Kosten für einen zusätzlichen Kubikmeter.“ Beurteile die Aussage \hfill (Abitur 2018 GK).
\end{teile}
\end{aufgabe}

%% TODO Titel: Ich-kann-Satz fehlt in der Bank (Platzhalter: Kettenname) – Nr. 11
%% TODO Anweisung: ein Satz über der ersten Teilaufgabe fehlt in der Bank – Nr. 11
\begin{aufgabe}{Mittlere Änderungsrate – Fehler finden, Begründen, Anwendung, Darstellungswechsel}
\begin{teile}
\teil Um 9:00 Uhr sind 250 Besucher im Museum, um 13:00 Uhr 850. Lea gibt als mittlere Änderungsrate an: „$850 - 250 = 600$ Besucher je Stunde.“ Finde den Fehler.
\teil Begründe, warum die mittlere Änderungsrate von $f$ auf $[a; b]$ genau die Steigung der Sekante durch $(a | f(a))$ und $(b | f(b))$ ist.
\teil Ein Wassertank fasst zu Beginn $800$ Liter, nach $6$ Stunden noch $380$ Liter. Der Betreiber darf höchstens $60$ Liter je Stunde verbrauchen. Hält er die Vorgabe im Mittel ein?
\teil Der Graph zeigt den Wasserstand $w$ (in cm) eines Bachs in Abhängigkeit von der Zeit $t$ (in h). Zeichne die Sekante durch die Punkte bei $t = 1$ und $t = 3$ ein und gib ihre Steigung als mittlere Änderungsrate mit Einheit an.

\begin{ksys}[xmin=0,xmax=4,ymin=0,ymax=8,xlabel=t in h,ylabel=w in cm,ablesen]\funktion{2+\x^2/2+\x/2}{w}\end{ksys}
\end{teile}
\end{aufgabe}
